
\subsubsection{6.1 Root Causes of Failure}

\textbf{Dataset-specific drift baselines:} The 18.8× variance in drift scores (0.042 to 0.79) across datasets with similar ground truth labels indicates that drift magnitude is relative to dataset-specific baselines rather than absolute. SST2 (labeled PATCH) scored 0.79 while GLUE QNLI (labeled MINOR) scored 0.042—a 18.8× difference. Both are GLUE benchmark train/validation splits with similar construction, yet drift scores differ by nearly two orders of magnitude.

This variance suggests that each dataset has an intrinsic baseline drift level. For SST2, a drift of 0.79 may represent normal train/validation variation, while for QNLI, the same score would represent a substantial shift. The fixed thresholds (0.07, 0.02, 0.005) cannot accommodate this variance.

\textbf{Cross-modality differences:} The thresholds (7\%/2\%/0.5\%) were derived from ImageNet literature. The tested datasets included 13 NLP examples and 1 vision example. NLP drift scores ranged from 0.042 to 0.79. The vision example (MNIST→USPS) scored 0.597, but this represents cross-dataset domain shift rather than version drift within a single dataset, limiting its validity for threshold calibration assessment.

\textbf{Feature extractor robustness:} The frozen pre-trained models (ResNet-50, BERT-base) were trained for transfer learning, which prioritizes robustness across distribution shifts. Models optimized to generalize well across domains may be insufficiently sensitive to detect subtle version-level changes. The 25\% recall suggests that frozen feature extractors may not capture version-relevant distribution differences.

\subsubsection{6.2 SST2 Anomaly}

SST2 showed the highest drift score (0.79) despite being labeled PATCH. Three possible explanations exist:

\begin{enumerate}
\item The PATCH label may be incorrect; SST2 validation may represent a more substantial shift than documented.
\item Validation set curation (e.g., class balancing, filtering) may introduce statistical drift without affecting model performance.
\item BERT embeddings for short sentences may be less stable than for longer documents.
\end{enumerate}

Without measured model performance (training on SST2-train and testing on SST2-validation), these hypotheses cannot be distinguished. If accuracy drop is <2\%, this confirms the PATCH label and reveals a false positive in drift detection. If accuracy drop is >5\%, the ground truth label should be revised.

\subsubsection{6.3 Limitations}

\textbf{Ground truth validation:} Labels were derived from literature descriptions rather than measured model performance. This means "ground truth" represents expert judgment about dataset construction rather than empirical performance impact. The 28.57\% accuracy could underestimate method performance if ground truth labels contain errors, or accurately reflect poor calibration if labels are correct.

\textbf{Dataset coverage:} Only 1 of 4 planned vision datasets was tested, and that dataset (MNIST→USPS) represents cross-dataset domain shift rather than version drift within MNIST. Claims about cross-modality generalization cannot be supported with this coverage. The tested sample consists primarily of NLP benchmarks (13 of 14 valid datasets).

Six datasets were unavailable: ImageNet-v2 and CIFAR-10.1 require manual download and registration; Fashion-MNIST variants encountered HuggingFace API errors; QQP, SQuAD, and MS-MARCO had API timeouts or configuration issues. The hypothesis failed on the available datasets with sufficient margin (-56.43pp accuracy gap) that the missing datasets are unlikely to reverse the conclusion.

\textbf{MNIST contamination:} MNIST→USPS/EMNIST tests cross-dataset domain shift (different digit sources) rather than version drift within MNIST, as no MNIST version updates exist. This data point is invalid for testing within-dataset version drift detection. Excluding it reduces the sample to 13 datasets but does not change the conclusion given the below-random accuracy.

\textbf{Feature extractor sensitivity:} Using frozen pre-trained models (optimized for robustness) may have reduced sensitivity to version-level shifts. However, even with current drift scores showing 18.8× variance, the core problem—threshold mis-calibration across datasets—persists. More sensitive features might improve recall but would not resolve the systematic precision issues arising from dataset-specific baselines.

\subsubsection{6.4 Implications}

This work provides evidence that drift magnitude alone, measured through standard statistical tests with fixed thresholds, is insufficient for semantic version classification without dataset-specific calibration. The 18.8× variance in drift scores across datasets with similar ground truth labels demonstrates that universal thresholds cannot exist in their tested form.

Three alternative approaches could address the identified failure modes:

\begin{enumerate}
\item \textbf{Adaptive per-dataset calibration:} Establish dataset-specific thresholds after observing multiple version transitions. This addresses the baseline variance problem but requires sufficient version history per dataset.
\end{enumerate}

\begin{enumerate}
\item \textbf{Supervised classification:} Train classifiers on labeled version pairs using features including drift score, dataset size, modality, and other metadata. This learns decision boundaries rather than assuming fixed thresholds but requires expensive labeled training data.
\end{enumerate}

\begin{enumerate}
\item \textbf{Performance-based ground truth:} Replace statistical drift with measured model performance degradation. Train models on v_old, measure accuracy on v_new, and use performance drops to define severity. This directly measures impact but requires computational resources and task-specific model architectures.
\end{enumerate}
